\documentclass[oneside]{scrbook}

% \usepackage[utf8]{inputenc}
\usepackage[T1]{fontenc}
\usepackage{version}
% \usepackage[latin1]{inputenc}
\usepackage[italian]{babel}
\usepackage{algpseudocode}
\usepackage{amsmath}
\usepackage{booktabs}
\usepackage{color}
\usepackage{amsfonts}
\usepackage{url}
\usepackage{algorithm}
\usepackage{chngcntr}
\usepackage[answerdelayed]{exercise}
\usepackage{graphicx}
\usepackage{caption}
\usepackage{listings}
\usepackage{tikz}
\usepackage{pgf}
\usepackage{colortbl}
\usepackage{listingsutf8}
\usepackage[pdftitle = {Dispense di esercizi}, pdfauthor = {Gianmarco Tedeschi e Alberto Marchesi}, pdfsubject = {Fondamenti di Informatica ($082746$ - A.A.~$2025-2026$)}, pdfcreator = {Gianmarco Tedeschi e Alberto Marchesi}, pdfproducer = {-}]{hyperref}
% \usepackage[automark,nouppercase]{scrpage2}
\usepackage{mdframed}
\usepackage{scrlayer-scrpage} % Necessary package for page styles

%\pagestyle{scrheadings}

%\usepackage{titlesec}
%\newcommand{\sectionbreak}{\clearpage}
\counterwithin{Exercise}{chapter}
\counterwithin{Answer}{chapter}

\algrenewcommand\algorithmicthen{\textbf{Allora}}
\algrenewcommand\algorithmicif{\textbf{Se}}
\algrenewcommand\algorithmicelse{\textbf{Altrimenti}}
\algrenewcommand\algorithmicend{\textbf{Chiudi}}
\algrenewcommand\algorithmicwhile{\textbf{Finch\'e}}
\algrenewcommand\algorithmicdo{\textbf{Esegui}}
\algrenewcommand{\algorithmiccomment}[1]{\hskip3em\texttt{//#1}}

\setlength{\parindent}{0in}
\setlength{\parskip}{10pt}

\newmdenv[%
        outerlinewidth=2,%
        roundcorner=10pt,%
        leftmargin=0,%
        rightmargin=0,%
        outerlinecolor=black,
        innertopmargin=\topskip,%
        splittopskip=\topskip,%
]{algo}


\ohead{\headmark}
\ihead{}
\chead{}
\ifoot{A.A.~$2025-2026$}
\cfoot{Fondamenti di Informatica (082746) - G. Tedeschi}
\ofoot{Pag.~\pagemark}
\setheadsepline{.4pt}
\setfootsepline{.4pt}

\definecolor{lbcolor}{rgb}{0.9,0.9,0.9}

\lstdefinestyle{cstyle}{%
  inputencoding=utf8/latin1,
  breaklines=true,
  breakatwhitespace=false,
  basicstyle=\small\ttfamily,
  language=c,
  frame=single,
  showstringspaces=false
}

\lstdefinestyle{cstylesmall}{%
  inputencoding=utf8/latin1,
  breaklines=true,
  breakatwhitespace=false,
  basicstyle=\scriptsize\ttfamily,
  language=c,
  frame=single,
  showstringspaces=false
}

\lstdefinestyle{console}{%
  inputencoding=utf8/latin1,
  breaklines=true,
  breakatwhitespace=false,
  basicstyle=\scriptsize\ttfamily,
  showstringspaces=false,
  backgroundcolor=\color{lbcolor}
}


\lstnewenvironment{clang}{\lstset{style=cstyle}}{}
\lstnewenvironment{clangsmall}{\lstset{style=cstylesmall}}{}

\addtocounter{chapter}{2}

\begin{document}

\chapter{Costrutti \texttt{while}, \texttt{for} e \texttt{switch}}

Questa dispensa propone esercizi sulla scrittura di algoritmi, in linguaggio C, utili alla comprensione dei costrutti \texttt{while}, \texttt{for} e \texttt{switch}.

I costrutti per costruire cicli in C sono il \texttt{while}, la variante \texttt{do...while} e il \texttt{for}.

\begin{lstlisting}[style=cstyle]
  inizializzazione; //opzionale

  while (condizione) {
    corpo;

    incremento; //opzionale
  }
\end{lstlisting}

La \texttt{condizione} \`e valutata prima di ogni iterazione, inclusa la prima; quindi il corpo del ciclo (e quindi anche l'eventuale istruzione di \texttt{incremento}) potrebbero non eseguire mai, nel caso in cui la condizione sia falsa dall'inizio.

\begin{lstlisting}[style=cstyle]
  inizializzazione; //opzionale

  do {
    corpo;

    incremento; //opzionale
  } while (condizione);
\end{lstlisting}

Invece, nella variante \texttt{do...while}, il corpo e l'eventuale istruzione di incremento sono eseguiti almeno una volta prima di valutare la condizione.

Il \texttt{for} è equivalente, ma ha una sintassi più compatta:

\begin{lstlisting}[style=cstyle]
  for (inizializzazione; condizione; incremento) {
    corpo;
  }
\end{lstlisting}

Si noti che le parentesi sono necessarie solo nel caso di corpo con più di un'istruzione. Inoltre, l'espressione di inizializzazione e di incremento sono opzionali. Di fatto, un ciclo \texttt{while} può essere scritto in modo equivalente con un costrutto \texttt{for} nel seguente modo:
\newpage
\begin{lstlisting}[style=cstyle]
  inizializzazione;

  for (; condizione; ) {
    corpo;

    incremento;
  }
\end{lstlisting}

Il costrutto del C che permette di scegliere tra alternative multiple (pi\`u di $2$) \`e lo \texttt{switch}.

\begin{lstlisting}[style=cstyle]
  inizializzazione;

  switch (variabile) {
    case valore1:
    	istruzione1;
    	break;
    case valore2:
    	istruzione2;
    	break;
    default:
    	istruzione3;
    	break;
  }
\end{lstlisting}

Il \texttt{break} tra un case e l'altro permette che venga eseguito solo il codice relativo al case in cui si \`e capitati. Il \texttt{break} finale non \`e necessario ma consigliato (nel caso in cui si dovesse immettere nuovi case).

\subsection{Esercizi}

\begin{Exercise}[label={ex:somma}]
  	Data la seguente operazione matematica:
  	$$ y = \sum^{n}_{i=0} i$$
    Si scriva un programma per il calcolo di $y$. Si assuma di acquisire $n$ come dati di input. Prima di procedere con il calcolo di $y$, si controlli che $n$ sia un intero positivo.
\end{Exercise}

\begin{Answer}[ref={ex:somma}]
  \lstinputlisting[style=cstyle]{src/somma.c}
\end{Answer}

\begin{Exercise}[label={ex:radice}]
  Scrivere un programma che dato un numero positivo ne restituisca la radice intera
\end{Exercise}

\begin{Answer}[ref={ex:radice}]
	Abbiamo già visto in un'esercitazione precedente lo pseudocodice per questo algoritmo.
	
	Iniziamo con la soluzione utilizzando il costrutto \texttt{while}
  \lstinputlisting[style=cstyle]{src/radice_intera1.c}

	Alternativamente si può utilizzare anche il ciclo \texttt{for}
	\lstinputlisting[style=cstyle]{src/radice_intera2.c}

	\paragraph{Attenzione: }
			In questa soluzione il ciclo non è vuoto, ma semplicemente l'unica operazione è il decremento della variabile che viene eseguito nella terza espressione del ciclo \texttt{for}
\end{Answer}


\begin{Exercise}[label={ex:switch_easy}]

    \Question Realizzare un programma che legga da input un voto (carattere da 'A' ed 'E') e ne stampi il significato (A significa ottimo, E significa gravemente insuff.,etc.)
    
    \Question Supporre che l'utente possa inserire valori definiti con una codifica errata. Si ripeta l'acquisizione dei dati finchè non vengono inseriti valori validi.

\end{Exercise}

\begin{Answer}[ref={ex:switch_easy}]
		\Question \lstinputlisting[style=cstyle]{src/switch_easy.c}

        \Question \lstinputlisting[style=cstyle]{src/switch_easy_2.c}

        \paragraph{Nota:} alternativamente, l'esercizio poteva essere risolto utilizzando una serie di costrutti \texttt{if} nidificati

        \lstinputlisting[style=cstyle]{src/switch_easy_3.c}
\end{Answer}
\begin{Exercise}[label={ex:aree}]
		\Question Scrivere un programma che dato un numero reale positivo $l$ e un intero positivo $n$ restituisca l'area del poligono regolare con $n$ lati di lunghezza $l$. Si implementi una soluzione per che gestisca i casi $n \in \{ 1, \ldots, 6 \}$ e che sia espandibile agevolmente.
		Suggerimento: l'area di un poligono regolare \`e:
		\begin{equation}
			A = \frac{nl^2}{4\tan(\frac{\pi}{n})}
		\end{equation}
		
		\Question Supporre ora che l'utente possa inserire valori negativi per $l$ e $n$. Si ripeta l'acquisizione dei dati finché l'utente non inserisce dei valori accettabili per i due parametri.
\end{Exercise}

\begin{Answer}[ref={ex:aree}]
  \Question \lstinputlisting[style=cstyle]{src/aree.c}

  \Question \lstinputlisting[style=cstyle]{src/aree2.c}
\end{Answer}


\begin{Exercise}[label={ex:divisori}]
  Scrivere un programma che richiede all'utente un intero positivo e ne stampa a schermo tutti i divisori.
\end{Exercise}

\begin{Answer}[ref={ex:divisori}]
  \lstinputlisting[style=cstyle]{src/divisori_while.c}

  \lstinputlisting[style=cstyle]{src/divisori_for.c}
  \paragraph{Attenzione: }
  		L'\verb+init_expr i = 1;+ e la \verb-loop_expr i++-; fanno parte della prima riga nel ciclo. Nel ciclo \verb+while+ erano prima del ciclo e all'interno rispettivamente.
\end{Answer}

\begin{Exercise}[label={ex:primo}]
  Scrivere un programma che richiede all'utente un intero positivo e determina se è primo o meno. Il programma deve continuare a chiedere il numero fino a che l'utente non ne inserisce uno positivo.
\end{Exercise}

\begin{Answer}[ref={ex:primo}]
  \lstinputlisting[style=cstyle]{src/primo.c}

  \lstinputlisting[style=cstyle]{src/primo-alternativo.c}
\end{Answer}
\begin{Exercise}[label={ex:tavola}]
  Scrivere un programma che stampi a video la tavola pitagorica da 1-$N$. Dove $N$ viene inserito dall'utente.
\end{Exercise}

\begin{Answer}[ref={ex:tavola}]
  \lstinputlisting[style=cstyle]{src/tavola.c}
\end{Answer}

\begin{Exercise}[label={ex:primi}]
  Scrivere un programma che richiede all'utente un intero positivo $N$ e stampa a schermo i primi $N$ numeri primi. Ad esempio: con $N = 7$ a schermo avremo $2$ $3$ $5$ $7$ $11$ $13$ $17$.
\end{Exercise}

\begin{Answer}[ref={ex:primi}]
  \lstinputlisting[style=cstyle]{src/primi.c}
\end{Answer}

% \begin{Exercise}[label={ex:segreto}]
  Scrivere un programma che richiede all'utente un indovinare un numero casuale (generato con la funzione \texttt{rand}) tra $1$ e $10$. Il gioco si svolge nel seguente modo: se l'utente ha indovinato allora viene stampato a schermo il successo, altrimenti il programma dovr\`a guidare l'utente dandogli come suggerimento se il numero segreto \`e pi\`u grande o pi\`u piccolo di quello inserito.
\end{Exercise}

\begin{Answer}[ref={ex:segreto}]
		\lstinputlisting[style=cstyle]{src/segreto.c}
\end{Answer}

\begin{Exercise}[label={ex:mcd}]
  Scrivere un programma che richiede due numeri e ne stampa il massimo comune divisore (MCD) e il minimo comune multiplo (mcm).
\end{Exercise}

\begin{Answer}[ref={ex:mcd}]
  \lstinputlisting[style=cstyle]{src/mcd.c}
  \paragraph{Nota: } Un metodo più efficiente per trovare MCD è l'utilizzare l'algoritmo di Euclide.
  \lstinputlisting[style=cstyle]{src/mcd_2.c}
\end{Answer}
\begin{Exercise}[label={ex:scf}]
  Implementare il gioco della morra cinese, noto anche come ``carta-forbice-sasso''. Il programma si aspetta da tastiera l'input dal giocatore 1 e, dopo, l'input dal giocatore 2. L'input potrà essere uno dei caratteri `C', `F' o `S'. Il programma dovrà continuare a chiedere input ai due giocatori fino a che ognuno dei due giocatori non inserisca una giocata valida. In questo caso calcola il vincitore (giocatore 1 o 2 o pareggio), stampa a video il vincitore e termina.
\end{Exercise}

\begin{Answer}[ref={ex:scf}]
  \lstinputlisting[style=cstyle]{src/csf.c}
\end{Answer}
\begin{Exercise}[label={ex:sequenza}]
  Realizzare un programma che prenda in input una sequenza di cifre (tra 1 e 9) e calcoli la somma massima fra le sotto-sequenze di cifre non decrescenti. Il programma termini quando viene inserito lo 0.
  
  Esempio:
  \begin{tabular}{|c c c c|c c|c|c c|c|}
    \hline
    2 & 2 & 4 & 5 & 3 & 9 & 3 & 1 & 5 & 0 \\
    \hline
    \textbf{13} & & & & 12 & & 3 & 6 & & \\
    \hline
  \end{tabular}

\end{Exercise}

\begin{Answer}[ref={ex:sequenza}]
  \lstinputlisting[style=cstyle]{src/sequenza.c}
\end{Answer}
\begin{Exercise}[label={ex:amicali}]
		Dati due numeri $m$ e $n$ questi sono numeri amicali se la somma dei divisori di $m$ è uguale a $n$, e viceversa (per esempio $220$ e $284$, $1184$ e $1210$, $2620$ e $2924$, $5020$ e $5564$, $6232$ e $6368$, $17296$ e $18416$). Ad esempio si ha:
		\begin{align*}
			& 220 \rightarrow 110 + 55 + 44 + 22 + 20 + 11 + 10 + 5 + 4 + 2 + 1 = 284\\
			& 284 \rightarrow 142 + 71 + 4 + 2 + 1 = 220
		\end{align*}
		Scrivere un programma che ricevuto in ingresso due numeri interi restituisce $1$ se i numeri sono amicali, $0$ altrimenti.
\end{Exercise}

\begin{Answer}[ref={ex:amicali}]
		\lstinputlisting[style=cstyle]{src/amicali.c}
		\paragraph{Attenzione: }
			L'algoritmo che viene presentato non \`e efficiente. Volendo si potrebbe minimizzare il numero delle iterazioni considerate.
\end{Answer}
\begin{Exercise}[label={ex:palindromo}]
  Scrivere un programma che richiede un numero e stampa se è palindromo o meno. Ad esempio il numero 121 è palindormo, mentre il numero 123 non lo è.
\end{Exercise}

\begin{Answer}[ref={ex:palindromo}]
  \lstinputlisting[style=cstyle]{src/palindromo.c}
\end{Answer}
\begin{Exercise}[label={ex:interessi}]
  Si progetti in C un programma che legga un float, rappresentante un ammontare di euro; di seguito il programma deve leggere un tasso d'interesse (in percentuale), ed un numero di anni. Il programma deve stampare, in uscita, per ogni anno, come l'ammontare cresce con gli interessi. Si ricordi che l'interesse si calcola con la seguente formula:
  $$
    C_{fin} = C_{in} * \left(1+\frac{r}{100}\right)^N
  $$
  Dove $C_{fin}$ è il capitale finale, $C_{in}$ è quello iniziale, $r$ è l'interesse e $N$ reppresenta il numero di anni in cui si applicano gli interessi.
\end{Exercise}

\begin{Answer}[ref={ex:interessi}]
  \lstinputlisting[style=cstyle]{src/interessi.c}
\end{Answer}




%TODO
\begin{Exercise}[label={ex:taylor}]
  	Scrivere un programma che dato un numero reale $x$ e un intero positivo $n$ calcoli lo sviluppo di McLaurin del seno in $x$ all'ordine $n$. Si ricorda che lo sviluppo di Taylor del seno \`e dato dalla sommatoria:
  	$$ \sin{x} = \sum^{n}_{i=0} \frac{(-1)^i}{(2i+1)!} x^{2i+1}$$
\end{Exercise}

\begin{Answer}[ref={ex:taylor}]
  \lstinputlisting[style=cstyle]{src/taylor.c}
  \paragraph{Nota: } Il ricalcolo del fattoriale ad ogni iterazione è costoso ed inefficiente, possiamo ottenere una soluzione più efficiente calcolandolo in modo incrementale.  
  \lstinputlisting[style=cstyle]{src/taylor_2.c}
\end{Answer}


\clearpage
\section*{Soluzioni}
\shipoutAnswer

\end{document}
